\reviewercomment{3. Classification of coalescence outcomes

The classification into Dominance, Mixture, and Restructuring is useful and provides a clear framework for summarising outcomes. However, it depends on thresholding a continuous similarity measure.

As noted in the broader literature [Mansour et al., 2018], coalescence outcomes often lie along a continuum rather than discrete categories. In addition, the distinction between mixture and restructuring is not entirely straightforward biologically, particularly because abundance changes contribute to classification.

We suggest that the authors:
- present continuous similarity measures alongside categorical outcomes
- clarify the biological interpretation of "restructuring"

We also find the comparison between dot product and Jaccard metrics in Extended Data Fig. 2 potentially very informative. The divergence between these metrics in low-nutrient environments may itself provide evidence that species-level processes dominate in that regime, and this could be discussed more explicitly.}

\responselabel We agree that the three outcome classes are a discretization of an underlying continuous similarity space, and we have revised the manuscript to make the continuous structure, the threshold dependence, and the biological meaning of Restructuring explicit.

In the revised Methods and Supplementary Methods, the classification is now described in terms of two continuous coordinates: retention magnitude $r = \sqrt{u^2+v^2}$, which measures how much of the coalesced community is explained by the two parental-community composition vectors, and PDI, which measures the angular bias toward one parental community versus the other. The categorical thresholds are now stated as interpretive boundaries in this space: Restructuring when $r^2 \leq 0.5$, Mixture when $r^2 > 0.5$ and $0.25 \leq \mathrm{PDI} \leq 0.75$, and community-level selection (CLS, called Dominance in the original classification figure) when $r^2 > 0.5$ and PDI is outside that interval. We also added continuous per-medium analyses in Supplementary Figs.~10--12 and a boundary-sensitivity analysis in Supplementary Fig.~37. Across the 263 synthetic coalescence events used in the nutrient-gradient analysis, the continuous asymmetry coordinate differs strongly across media (Kruskal--Wallis $H = 110.71$, $p = 9.12 \times 10^{-25}$), and the same nutrient-dependent shift is visible before applying category labels: median asymmetry increases from 0.410 in Nutr$-$ to 0.915 in Base medium and 0.985 in Nutr$+$. Retention magnitude also differs across media (Kruskal--Wallis $H = 18.05$, $p = 1.21 \times 10^{-4}$), with median $r^2$ values of 0.792, 0.722, and 0.924 in Nutr$-$, Base medium, and Nutr$+$, respectively. Thus, the categorical result is not only a threshold artifact: the underlying continuous coordinates shift systematically with nutrient condition.

\begin{figure}[H]
\centering
\includegraphics[width=0.92\textwidth]{figures/r2_q3_continuous_similarity_and_metric_divergence.png}
\caption{\textbf{Response Fig.~R2Q3. Continuous similarity coordinates and metric dependence.} \textbf{A,} Each synthetic coalescence event is plotted by squared retention magnitude $r^2$ and parental asymmetry. Dashed lines show the manuscript thresholds. \textbf{B,} Outcome fractions computed with the manuscript dot-product/vector-decomposition framework and with Jaccard presence/absence similarity. Jaccard shifts many low-nutrient events from CLS to Mixture, consistent with species-level persistence from both parental communities when interactions are weak. In Base medium and Nutr$+$, the Jaccard divergence mainly shows that abundance-weighted CLS does not require retaining the full species list of the winning parental community.}
\end{figure}

We have also clarified the biological interpretation of Restructuring. In the revised text, Restructuring does not mean a balanced combination of both parental communities. It means low parental retention: the best linear combination of the two parental-community profiles explains less than half of the coalesced composition in the normalized similarity space. Biologically, this can occur when cross-community interactions after mixing produce a new stable composition, for example through competitive exclusion cascades or growth of taxa that were not dominant in either parental community before coalescence. This interpretation is now stated in the Results and in the Supplementary Methods, and the response is bounded accordingly: Restructuring identifies low retention of the two parental-community profiles, but it does not by itself identify the causal interaction mechanism.

Finally, we agree that the divergence between dot product and Jaccard is biologically informative rather than only a technical robustness issue. We now describe dot product/vector decomposition as an abundance-weighted readout of parental-community dominance, whereas Jaccard is a presence/absence readout of species-identity retention. In our response-only recalculation across media, the Nutr$-$ dot-product classification gives 35 CLS, 48 Mixture, and 7 Restructuring events, whereas the Jaccard classification gives 5 CLS, 55 Mixture, and 30 Restructuring events. This supports the reviewer's interpretation that weak-interaction, low-nutrient coalescence often retains species from both parental communities even when abundance-weighted similarity is uneven. We use more cautious language for Base medium and Nutr$+$ because the same Jaccard divergence there is not evidence for simple species-level mixing; rather, it shows that abundance dominance can coexist with loss or rearrangement of the winning parental community's presence/absence profile.

\mschange{We revised the Results to define Restructuring as low parental retention, distinguish it from Mixture, and point readers to the new continuous distributions in Supplementary Figs.~10--12. We revised the Supplementary Methods to state the continuous coordinates, threshold rationale, and boundary-sensitivity analysis, added Supplementary Figs.~10--12 and Supplementary Fig.~37, and revised the Extended Data Fig.~2 caption and Similarity Metric Robustness paragraph to interpret Jaccard as a complementary species-identity retention readout rather than an abundance-dominance metric.}
